% !PRINTED 515
\weberpage{515}

The product (4) now admits the substitutions of the group $Q$,
and hence, if $\psi$, as above, denotes a function in $N$
belonging to this group, then by \S.\,155, 2. the function
$g_1(t,\varepsilon)$ can be expressed rationally in terms of
$\psi$; that is, $g_1(t,\varepsilon)$ is a function of $t$ in
the field of degree $j$, $\Omega(\psi)$, and shall accordingly
be denoted by $g(t,\psi)$.

The quantity $\psi$ is one of the roots of an irreducible
equation of degree $j$

\weq{6}{\varphi(u) = 0,}

whose remaining roots are $\psi_1, \psi_2 \dots \psi_{j-1}$,
and by \S.\,155, (6) $g(t)$ is decomposable into the factors

\weq{7}{g(t) = g(t,\psi)\, g(t,\psi_1) \dots g(t,\psi_{j-1}).}

Now

\[ g(t,\psi) = g_1(t,\varepsilon), \]

and from (7) it follows that this equation does not continue to
hold if, on the left-hand side, one of the other roots of (6) is
put in place of $\psi$. One can therefore assign to $t$ a
rational value such that the equation

\weq{8}{g(t,u) - g_1(t,\varepsilon) = 0}

has only the one root $u = \psi$ in common with equation (6).
The greatest common divisor of (6) and (8) is then linear with
respect to $u$, and if one sets $u = \psi$ in it, it results
\gesperrt{that $\psi$ is rationally expressible in terms of
$\varepsilon$}. The field $\Omega(\psi)$ is therefore certainly a
divisor of the field $\Omega(\varepsilon)$. The degree of the
field $\Omega(\varepsilon)$, i.\,e. the degree of the auxiliary
equation $\chi(\varepsilon) = 0$, is therefore a multiple of $j$
and never lower than $j$. If the degree of the auxiliary
equation is equal to $j$, which, for example, always occurs when
the degree of $\chi$ is a prime, then $\Omega(\varepsilon)$ is
identical with $\Omega(\psi)$; that is, $\varepsilon$ is also
rationally expressible in terms of $\psi$. In this case the
conjugate values $\varepsilon, \varepsilon_1 \dots
\varepsilon_{j-1}$ are all the roots of $\chi = 0$, and one
obtains them by replacing, in the rational expression of
$\varepsilon$ in terms of $\psi$, the function $\psi$ by each of
the conjugate functions $\psi, \psi_1 \dots \psi_{j-1}$. Hence
$\varepsilon$ itself can be taken for $\psi$, and one finds a
decomposition

\[ g(t) = g(t,\varepsilon)\, g(t,\varepsilon_1) \dots g(t,\varepsilon_{j-1}). \]

The quantities $\varepsilon, \varepsilon_1, \varepsilon_2 \dots
\varepsilon_{j-1}$ all belong to the field $\Omega(\varrho)$.
